\section{Data Card and Benchmark Documentation}
\label{app:datacard}

\subsection{Task Suite Descriptions}

\textbf{Synthetic trap tasks (95 tasks).} Programmatically generated across 5 domains (tabular EDA, time-series, recommendation, statistical, text-log) with deterministic ground-truth answers. Each task includes designed failure traps (e.g., count-vs-sum ambiguity, temporal leakage, type confusion). Data is generated at runtime from fixed seeds, ensuring reproducibility without external data dependencies.

\textbf{UCI real-data tasks (23 tasks).} Analytical questions on 6 public datasets (Iris, Titanic, Tips, MPG, Penguins, Flights) downloaded at runtime from public repositories. Answers are computed deterministically from the real data. Tasks cover groupby aggregation, correlation, counting, and filtering.

\textbf{DSBench real tasks (200 tasks).} Real financial-analysis questions from ModelOff competitions~\citep{jing2024dsbench}, using original multi-sheet Excel workbooks. Answers are multiple-choice (A--I) or numeric. Tasks require understanding complex financial data structures.

\subsection{Task Difficulty Distribution}
\label{app:difficulty}

We characterise task difficulty by the baseline success rate aggregated across all five models. \Cref{tab:difficulty} reports the distribution. Synthetic tasks span the full difficulty range (19 easy, 63 medium, 13 hard); UCI tasks are uniformly medium-hard; DSBench tasks are almost entirely hard (197 of 200 tasks have SR $<$ 25\%), reflecting the complexity of real financial data.

\begin{table}[t]
\centering
\caption{Task difficulty distribution by suite. Easy: SR $\geq$ 75\%; Medium: 25\% $\leq$ SR $<$ 75\%; Hard: SR $<$ 25\%.}
\label{tab:difficulty}
\small
\begin{tabular}{@{}lcccc@{}}
\toprule
Suite & Tasks & Easy & Medium & Hard \\
\midrule
Synthetic & 95 & 19 & 63 & 13 \\
UCI & 23 & 0 & 21 & 2 \\
DSBench & 200 & 0 & 3 & 197 \\
\bottomrule
\end{tabular}
\end{table}

\subsection{Expected Baseline Performance}
\label{app:baselines}

\Cref{tab:baselines} reports the expected baseline performance for each model and suite, providing reference points for future work. DeepSeek-V3 is the strongest on synthetic (88.4\%) and UCI (94.2\%) tasks. DSBench is challenging for all evaluated models (0--3.3\%).\begin{table}[t]
\centering
\caption{Baseline success rates (\%) by model and suite. ``n/a'' indicates API rate limiting prevented evaluation.}
\label{tab:baselines}
\small
\begin{tabular}{@{}lccc@{}}
\toprule
Model & Synthetic & UCI & DSBench \\
\midrule
DeepSeek-V3 & \textbf{88.4} & \textbf{94.2} & 0.0 \\
GPT-4o-mini & 73.9 & 82.6 & 0.8 \\
GPT-4o & 62.1 & 91.3 & \textbf{3.3} \\
Qwen3-max & 18.2 & 0.0 & 3.0 \\
DeepSeek-R1 & 0.0 & 0.0 & n/a \\
\bottomrule
\end{tabular}
\end{table}

\subsection{Data Schema}

Each run in the manifest (\texttt{manifest.csv}) contains the following fields:

\begin{itemize}[leftmargin=*,itemsep=1pt]
\item \texttt{task\_id}: unique task identifier (e.g., \texttt{tabular\_eda\_00})
\item \texttt{suite}: \texttt{synthetic}, \texttt{uci}, or \texttt{dsbench}
\item \texttt{model}: LLM identifier (e.g., \texttt{gpt-4o}, \texttt{deepseek-v3})
\item \texttt{variant}: \texttt{baseline} or \texttt{verifier}
\item \texttt{rep}: repetition index (0--2)
\item \texttt{correct}: boolean, whether the final answer matches ground truth
\item \texttt{failure\_stage}: \texttt{none}, \texttt{runtime}, or \texttt{silent}
\item \texttt{is\_silent}: boolean, whether the failure is silent
\item \texttt{tokens}: total tokens consumed
\item \texttt{final\_answer}: agent's final answer (truncated to 200 chars)
\item \texttt{causal\_attributed}: boolean, whether oracle repair succeeded
\item \texttt{source\_file}: original detail JSON filename
\end{itemize}

\subsection{Provenance}

All experiments use deterministic data generation (fixed seeds for synthetic tasks) and temperature-0 LLM calls. The framework includes a MockLLM backend that reproduces the full pipeline without API cost. Real-model experiments use the ChatAnywhere OpenAI-compatible endpoint. Total API cost: \$80.

\subsection{Licensing and Maintenance}

The code is released under an open-source license (MIT). The synthetic task generators are self-contained and require no external data. UCI datasets are public domain. DSBench tasks are used under their original license~\citep{jing2024dsbench}; users must obtain DSBench separately. The benchmark can be extended with new tasks by implementing the \texttt{Task} interface.

\textbf{Maintenance plan.} The repository is maintained by the authors. Issues and pull requests are accepted via GitHub. The maintainers commit to:
\begin{itemize}[leftmargin=*,itemsep=1pt]
\item Responding to issues within 30 days.
\item Releasing compatibility patches for new LLM API versions.
\item Preserving backward compatibility of the manifest schema across minor versions.
\item Archiving major versions as tagged releases.
\end{itemize}

\textbf{Versioning.} The benchmark follows semantic versioning (\texttt{MAJOR.MINOR.PATCH}). The current release is v1.0.0. Breaking changes to the manifest schema or task definitions increment \texttt{MAJOR}; new tasks or metrics increment \texttt{MINOR}; bug fixes increment \texttt{PATCH}. Each release is tagged on GitHub with a corresponding Zenodo DOI for long-term archival.

\textbf{Intended use cases.} The benchmark is designed for: (1) evaluating LLM-agent failure modes on data-science tasks; (2) comparing diagnostic metrics (SFR, propagation depth, recovery rate) across models; (3) testing new verifier or repair strategies. It is \emph{not} intended for: (1) training LLMs (tasks have deterministic answers); (2) production deployment evaluation (the minimal ReAct harness differs from production systems); (3) cross-domain generalisation claims (the three suites cover data analysis, not all of data science).

\subsection{Propagation Depth Analysis}

\Cref{tab:propagation} reports the propagation depth distribution across all 3{,}940 failures. The average depth is 0.30 steps, meaning most failures are detected immediately. However, 262 failures (6.9\%) propagate to depth $\geq$ 2, indicating long-horizon failure spread where the agent continues building on a wrong intermediate result.

\begin{table}[t]
\centering
\caption{Propagation depth distribution (all 3{,}940 failures).}
\label{tab:propagation}
\small
\begin{tabular}{@{}ccc@{}}
\toprule
Depth & Failures & Percentage \\
\midrule
0 & 3{,}350 & 88.6\% \\
1 & 169 & 4.5\% \\
2 & 80 & 2.1\% \\
3 & 39 & 1.0\% \\
4 & 64 & 1.7\% \\
5 & 45 & 1.2\% \\
6 & 20 & 0.5\% \\
7 & 14 & 0.4\% \\
\midrule
$\geq$ 2 (long-horizon) & 262 & 6.9\% \\
\bottomrule
\end{tabular}
\end{table}
